% !TeX root = ../tesis.tex

\chapter*{Introducción}
\addcontentsline{toc}{chapter}{\protect\numberline{}Introducción}	  		% Comment if you don't want the introduction to appear on the table of content. It will not have a number
\label{chapter:intro}


El esparcimiento de luz por partículas depende de propiedades como el tamaño, la geometría, la composición, la estructura interna y la orientación del sistema considerado~\cite{baldiniShiningLightFuture2025}. Esta dependencia hace que la respuesta óptica sea sensible a cambios en las propiedades físicas y composicionales de los sistemas biológicos, lo que permite emplearla como una fuente de información para su caracterización. Por esta razón, la interacción de la luz con la materia biológica se ha aprovechado en el desarrollo de métodos de diagnóstico, monitoreo y tratamiento~\cite{baldiniShiningLightFuture2025}. En particular, las alteraciones de las células y tejidos pueden modificar su interacción con el campo electromagnético incidente, generando cambios en su respuesta óptica que pueden emplearse para su caracterización~\cite{baldiniShiningLightFuture2025}.

Un ejemplo de ello es el estudio del esparcimiento de luz en la sangre, cuya respuesta óptica puede mostrar cambios en las propiedades morfológicas y composicionales de sus componentes celulares asociados a distintas enfermedades \cite{ergulComputationalStudyScattering2010}. Entre los principales componentes celulares de la sangre se encuentran los eritrocitos, cuya función fisiológica fundamental es el transporte de oxígeno \cite{palTextbookMedicalPhysiology2021}. En el rango visible, los eritrocitos, constituidos principalmente por hemoglobina,  dominan la respuesta electromagnética de la sangre \cite{bosschaartLiteratureReviewNovel2014}. En condiciones fisiológicas, los eritrocitos presentan una geometría discoidal bicóncava, sin embargo, distintas alteraciones hematológicas pueden modificar tanto su morfología como su composición, produciendo variaciones en parámetros como la forma, el volumen y la concentración de hemoglobina corpuscular media (CHCM)~\cite{palTextbookMedicalPhysiology2021}, como se muestra en la Fig. \ref{subfig:eritro} donde se observa una micrografía de un eritrocito sano y un tipo de eritrocito enfermo conocido como esferocito asociado a la enfermedad esferocitosis hereditaria.

Estas modificaciones dan lugar a respuestas espectrales y patrones de esparcimiento diferentes, que pueden emplearse para caracterizar enfermedades, particularmente en el espectro visible, donde la hemoglobina presenta una estructura espectral característica, dominada por la banda de Soret y las bandas $\alpha$ y $\beta$, asociadas a transiciones electrónicas de la molécula \cite{dayerBandAssignmentHemoglobin2010}. Estas bandas producen variaciones en las partes real e imaginaria del índice de refracción y, por tanto, en los procesos de absorción y esparcimiento. Por ello, cambios en la morfología y en la CHCM pueden modificar tanto la respuesta espectral como la distribución angular de la luz esparcida por los eritrocitos \cite{bectondickinsonIntroductionFlowCytometry2000}. 

El cambio de la respuesta óptica asociado a estas propiedades ha motivado el desarrollo de técnicas ópticas para la identificación y diferenciación celular. Entre ellas, la citometría de flujo, que permite analizar y diferenciar grandes poblaciones de células [ver Fig. \ref{subfig:Flow_cyt1}] a partir de la luz esparcida al atravesar un láser [ver Fig. \ref{subfig:Flow_cyt2}] \cite{bectondickinsonIntroductionFlowCytometry2000}. No obstante, los detectores de los citómetros convencionales registran el esparcimiento únicamente dentro de determinadas regiones angulares, definidas por la geometría del sistema de medición \cite{mckinnonFlowCytometryOverview2018}. Por ello, conocer la distribución angular del esparcimiento en regiones angulares más amplias puede aportar información adicional para identificar regiones de observación en las que las diferencias entre distintas morfologías celulares sean más pronunciadas.
%
\begin{figure}[tb]
	\centering %

	\sidesubfloat[First image]{\hspace{-0.2cm}{
			\includegraphics[width=0.35\textwidth]{0-Introduccion/eritrocitos.png}\label{subfig:eritro}}}\hspace{0.5cm}
	\sidesubfloat[Second image]{\hspace{-0.2cm}{	\includegraphics[width=0.4\textwidth]{0-Introduccion/flow_citometry_1.png}\label{subfig:Flow_cyt1}}}
		\vspace{0.3cm}
	\sidesubfloat[Third image]{\hspace{-0.2cm}{	\includegraphics[width=0.83\textwidth]{0-Introduccion/flow_citometry_0.png}\label{subfig:Flow_cyt2}}}
	\caption{\textbf{a)} Eritrocitos observados mediante microscopía electrónica de barrido (SEM, por sus siglas en inglés). Imagen extraída de la Ref. \cite{vivescorronsRareAnemiasAdolescents2021} \textbf{b)} Ejemplo de un diagrama FSC vs SCC (Sección transversal de esparcimiento en la dirección lateral en función de la sección transversal de esparcimiento en la dirección frontal ) de una muestra de sangre. Las tres poblaciones principales ---linfocitos, monocitos y granulocitos--- pueden distinguirse según su tamaño relativo y su granularidad. Imagen extraída de la Ref. \cite{proserpioFlowCytometryBeginners2022}. \textbf{c)} Esquema de citometría de flujo. Un láser incide en una célula y el esparcimiento de luz por la célula es medido por detectores en la dirección frontal y lateral. Para células marcadas con fluorocromos, se emplean detectores de fluoresencia.  Imagen extraída de la Ref. \cite{proserpioFlowCytometryBeginners2022} }
	\label{fig:Flow_cyt}
\end{figure}
%

 Debido a la geometría bicóncava de los eritrocitos, el estudio del esparcimiento de luz por estás células se ha abordado mediante distintos métodos numéricos. Lu et al. \cite{luSimulationsLightScattering2005}, por ejemplo, emplearon el método de \textit{Finite Difference Time Domain} (FDTD) para calcular eficiencias de absorción y esparcimiento para un eritrocito sano a diferentes ángulos de incidencia del campo electromagnético [ver Fig. \ref{subfig:sca_pattern1}]; Eremina et al. \cite{ereminaDifferentShapeModels2006} emplearon el método de \textit{Discrete Sources Method} (DSM) para evaluar el patrón de esparcimiento de diferentes modelos de eritrocitos [ver Fig. \ref{subfig:sca_pattern2}] y Tsinopoulos et al. \cite{tsinopoulosLightScatteringAggregated2002} desarrollaron un enfoque basado en el \textit{Boundary Element Method} (BEM) para evaluar el efecto de la membrana de los eritrocitos en sus patrones de esparcimiento [ver Fig. \ref{subfig:sca_pattern3}]. No obstante, los trabajos mencionados anteriormente se concentran principalmente en longitudes de onda cercanas a 633 nm, fuera de las principales bandas de absorción de la hemoglobina. Esto limita el estudio de la dependencia espectral de la respuesta óptica en regiones donde las propiedades ópticas de los eritrocitos presentan variaciones más pronunciadas. Por lo tanto, resulta de interés extender el análisis a un intervalo espectral que incluya la banda de Soret y las bandas $\beta$ y $\alpha$, considerando además las modificaciones morfológicas y composicionales asociadas a eritrocitos enfermos.
%
\begin{figure}[tb]
	\centering %
	
	\sidesubfloat[First image]{\hspace{-0.2cm}{
			\includegraphics[width=0.45\textwidth]{0-Introduccion/QsQa.png}\label{subfig:sca_pattern1}}}\hspace{0.5cm}
	\sidesubfloat[Second image]{\hspace{-0.2cm}{	\includegraphics[width=0.44\textwidth]{0-Introduccion/Scattering_pattern.png}\label{subfig:sca_pattern2}}}
	\vspace{0.3cm}
	\sidesubfloat[Third image]{\hspace{-0.2cm}{	\includegraphics[width=0.44\textwidth]{0-Introduccion/Scattering_pattern_1.png}\label{subfig:sca_pattern3}}}
	\caption{\textbf{a)} Eficiencias de absorción y esparcimiento de un eritrocito sano a diferentes ángulos de incidente del campo electromagnético a una longitud de onda de 1000 nm. Imagen extraída de la Ref. \cite{luSimulationsLightScattering2005}. \textbf{b)} Patrón de esparcimiento de diferentes modelos de eritrocitos. Imagen extraída de la Ref. \cite{ereminaDifferentShapeModels2006}. \textbf{c)} Patrón de esparcimiento de un eritrocito sano con y sin membrana. Imagen extraída de la Ref. \cite{tsinopoulosLightScatteringAggregated2002} }
	\label{fig:sca_pattern}
\end{figure}
%


A partir de lo anterior, el objetivo de esta tesis es identificar condiciones espectrales y angulares en las que que aumente el contraste entre células sanas y enfermas. Para ello, se estudian tres morfologías representativas: un eritrocito sano, un esferocito y un microcito, asignando a cada una de ellas una CHCM y una morfología característica. Las propiedades ópticas se describen mediante funciones dieléctricas obtenidas a partir de relaciones de Kramers-Kronig sustractivas, mientras que su morfología se describe mediante óvalos de Cassini. La respuesta electromagnética se calcula mediante el método de \textit{Finite Integration Technique} (FIT) y los parámetros óptimos de simulación se establecen a partir de un análisis de convergencia empleando como referencia la solución analítica de Mie para una esfera y, posteriormente, el modelo de eritrocito se contrasta con resultados numéricos reportados en la literatura. Una vez establecidas las condiciones de simulación, se calculan las eficiencias de extinción, absorción y esparcimiento, así como se analizan los patrones de esparcimiento de las bandas de Soret, $\beta$ y $\alpha$ con el propósito de localizar intervalos de detección en los que aumente el contraste entre las distintas morfologías.

Esta tesis está dividida en cuatro capítulos:  en \hyperref[chapter:theory]{Esparcimiento de luz por partículas} se presenta el formalismo electromagnético empleado para describir el esparcimiento de luz por partículas, incluyendo las ecuaciones de Maxwell, las secciones transversales y las relaciones de Kramers-Kronig sustractivas. En \hyperref[chapter:eritrocitos]{Eritrocitos y sus patologías} se describen la morfología y composición de los eritrocitos, sus principales alteraciones, los fundamentos ópticos de la citometría de flujo y la construcción de las funciones dieléctricas empleadas en las simulaciones. En \hyperref[chapter:FIT]{\textit{Finite Integration Technique}} se introduce el método numérico empleado y se establece la configuración numérica mediante un análisis de convergencia. En \hyperref[chapter:Resultados]{Propiedades ópticas} se presentan y discuten los resultados correspondientes al eritrocito sano, el esferocito y el microcito, incluyendo su respuesta espectral, sus patrones angulares de esparcimiento y la identificación de regiones de mayor contraste. Finalmente, se presentan las conclusiones generales y las principales líneas de trabajo a futuro.